%% ============================================================
%% Introduction to Cognitive Psychology — Week 3 Study Notes
%% ============================================================
\documentclass[11pt,a4paper]{article}

\usepackage[a4paper, left=1.8cm, right=1.8cm, top=1.3cm, bottom=1.8cm]{geometry}
\usepackage{booktabs}
\usepackage{tabularx}
\usepackage{array}
\usepackage{enumitem}
\usepackage{titlesec}
\usepackage{fancyhdr}
\usepackage{fontenc}
\usepackage{inputenc}
\usepackage{microtype}
\usepackage{parskip}
\usepackage{hyperref}
\usepackage{amsmath}
\usepackage{amssymb}

\hypersetup{hidelinks, colorlinks=false}

%% ── Table helpers ─────────────────────────────────────────
\newcolumntype{L}[1]{>{\raggedright\arraybackslash}p{#1}}

%% ── Section headings ──────────────────────────────────────
\titleformat{\section}{\large\bfseries}{}{0em}{}[\titlerule]
\titleformat{\subsection}{\normalsize\bfseries}{}{0em}{}
\titleformat{\subsubsection}{\normalsize\bfseries}{}{0em}{}

\titlespacing{\section}{0pt}{12pt}{4pt}
\titlespacing{\subsection}{0pt}{8pt}{3pt}
\titlespacing{\subsubsection}{0pt}{6pt}{2pt}

%% ── Page style ────────────────────────────────────────────
\pagestyle{fancy}
\fancyhf{}
\renewcommand{\headrulewidth}{0pt}
\fancyfoot[C]{\small Cognitive Psychology --- Week 3 \textbar\ Page \thepage}

%% ── Misc helpers ──────────────────────────────────────────
\newcommand{\bb}[1]{\textbf{#1}}
\setlength{\parskip}{4pt}

%% ═══════════════════════════════════════════════════════════
\begin{document}
\setlength{\arrayrulewidth}{0.4pt}

\begin{center}
\bfseries\LARGE Introduction to Cognitive Psychology\par
\vspace{0.2cm}
\large Assignment 3 Detailed Solutions
\end{center}
\vspace{0.5cm}

%% ════════════════════════════════════════════════════════════
\section*{ASSIGNMENT 3 --- Full Solutions with Explanations}

\subsubsection*{Question 1}
\textit{According to Attenuation theory of Treisman, people tend to process:}
\medskip

\begin{tabular}{L{0.3cm} L{14.9cm}}
& A.\ Only to the level of physical characteristics \\
& B.\ Only to the level of linguistic characteristics, separating it into words \\
& C.\ At a semantic level, analysing for meaning most of the time \\
& \textbf{D.\ Only as much as is necessary to separate the attended from the unattended message \checkmark} \\
\end{tabular}

\vspace{0.3cm}
\noindent\textbf{Ans:} D — Only as much as is necessary to separate the attended from the unattended message\\[0.2cm]
\textbf{Why this answer:} \\
Treisman's \textbf{Attenuation Theory} proposes that unattended information is not completely blocked (as Broadbent's filter theory suggested) but rather \textbf{weakened (attenuated)}. The degree to which unattended information is processed depends on the demands of the task — people process only as much as is necessary to distinguish the attended message from the unattended one. If the unattended message contains information with a low enough activation threshold (e.g., one's own name), it can still be detected.
\vspace{0.4cm}

\subsubsection*{Question 2}
\textit{The "cocktail party effect" refers to the fact that shadowing performance is disrupted when \_\_\_ is embedded in the unattended message.}
\medskip

\begin{tabular}{L{0.3cm} L{14.9cm}}
& A.\ Backward speech \\
& \textbf{B.\ The listener's name \checkmark} \\
& C.\ A section of repeated words \\
& D.\ Music \\
\end{tabular}

\vspace{0.3cm}
\noindent\textbf{Ans:} B — The listener's name\\[0.2cm]
\textbf{Why this answer:} \\
The \textbf{cocktail party effect} demonstrates that personally meaningful stimuli — most notably \textbf{one's own name} — can break through selective attention even when presented in an unattended channel. This phenomenon was famously studied by \textbf{Moray (1959)}, who found that participants performing a shadowing task noticed their own name in the unattended ear approximately one-third of the time. It challenges strict filter theories by showing that some semantic processing of unattended information must occur.
\vspace{0.4cm}

\subsubsection*{Question 3}
\textit{In Kahneman's model of attention, allocation of mental resources is affected by preferences for certain kinds of tasks over others. These preferences are known as:}
\medskip

\begin{tabular}{L{0.3cm} L{14.9cm}}
& \textbf{A.\ Enduring dispositions \checkmark} \\
& B.\ Arousal states \\
& C.\ Late selected preferences \\
& D.\ Momentary intentions \\
\end{tabular}

\vspace{0.3cm}
\noindent\textbf{Ans:} A — Enduring dispositions\\[0.2cm]
\textbf{Why this answer:} \\
In \textbf{Kahneman's Capacity Model of Attention} (1973), attentional resources are allocated based on several factors. \textbf{Enduring dispositions} are long-term, stable attentional biases that reflect an individual's habitual preferences for certain types of tasks or stimuli. These are distinguished from \textbf{momentary intentions}, which are temporary, goal-directed decisions about where to focus attention. Together, enduring dispositions and momentary intentions influence how a limited pool of mental effort is distributed across competing tasks.
\vspace{0.4cm}

\subsubsection*{Question 4}
\textit{The word "cat" is \_\_ by the phrase "The dog chased the…." That is, the word cat is especially ready to be recognized or attended to.}
\medskip

\begin{tabular}{L{0.3cm} L{14.9cm}}
& A.\ Filtered \\
& B.\ Attenuated \\
& \textbf{C.\ Primed \checkmark} \\
& D.\ Suggested \\
\end{tabular}

\vspace{0.3cm}
\noindent\textbf{Ans:} C — Primed\\[0.2cm]
\textbf{Why this answer:} \\
\textbf{Priming} occurs when prior contextual information activates related concepts in memory, increasing the likelihood and speed of recognizing an expected stimulus. In this example, the phrase \textit{"The dog chased the…"} creates a strong semantic context that \textbf{primes} the word "cat" — making it especially ready to be recognized, attended to, or completed. Priming can be both \textbf{semantic} (meaning-based) and \textbf{repetition-based} (prior exposure).
\vspace{0.4cm}

\subsubsection*{Question 5}
\textit{When listening to a conversation, your attention is momentarily diverted when you hear your name spoken in a different conversation across the room. This is an example of the:}
\medskip

\begin{tabular}{L{0.3cm} L{14.9cm}}
& A.\ Filter effect \\
& B.\ Dichotic listening phenomenon \\
& \textbf{C.\ Cocktail party effect \checkmark} \\
& D.\ Attenuation effect \\
\end{tabular}

\vspace{0.3cm}
\noindent\textbf{Ans:} C — Cocktail party effect\\[0.2cm]
\textbf{Why this answer:} \\
This is a classic real-world example of the \textbf{cocktail party effect}. Even while selectively attending to one conversation, highly salient or personally relevant stimuli — such as hearing your own name — in an unattended stream can capture your attention. This demonstrates that unattended information is not entirely ignored; some degree of \textbf{semantic processing} must occur for such personally meaningful stimuli to be detected and to redirect attention.
\vspace{0.4cm}

\subsubsection*{Question 6}
\textit{Which of the following would NOT be a reasonable basis for filtering, according to Broadbent's model?}
\medskip

\begin{tabular}{L{0.3cm} L{14.9cm}}
& A.\ Whether the message was coming from your right or your left side \\
& B.\ The pitch of the voice reading the message \\
& C.\ The loudness of the voice reading the message \\
& \textbf{D.\ The language that the message was being read in \checkmark} \\
\end{tabular}

\vspace{0.3cm}
\noindent\textbf{Ans:} D — The language that the message was being read in\\[0.2cm]
\textbf{Why this answer:} \\
\textbf{Broadbent's Filter Theory} (1958) proposes that selective attention operates as an early filter that selects information based solely on \textbf{physical characteristics} — such as spatial location (ear), pitch, loudness, or voice quality. The filter acts \textit{before} semantic analysis, meaning the meaning or language of the message cannot serve as a basis for filtering. Language identification requires \textbf{semantic processing}, which occurs \textit{after} the filter in Broadbent's model.
\vspace{0.4cm}
\newpage

\subsubsection*{Question 7}
\textit{Cognitive psychologists use all of the following as criteria for determining whether an activity is automatic EXCEPT:}
\medskip

\begin{tabular}{L{0.3cm} L{14.9cm}}
& A.\ Whether it occurs intentionally \\
& B.\ Whether it gives rise to conscious awareness \\
& C.\ Whether it interferes with other activities \\
& \textbf{D.\ Whether it requires mental filtering \checkmark} \\
\end{tabular}

\vspace{0.3cm}
\noindent\textbf{Ans:} D — Whether it requires mental filtering\\[0.2cm]
\textbf{Why this answer:} \\
\textbf{Automatic processes} are defined by three key criteria: (1) they occur \textbf{without intention} (involuntarily), (2) they happen \textbf{without conscious awareness}, and (3) they produce \textbf{little or no interference} with other concurrent activities. "Mental filtering" is not one of the standard criteria used to distinguish automatic from controlled processes. Controlled (effortful) processes, by contrast, are intentional, require awareness, and consume attentional resources.
\vspace{0.4cm}

\subsubsection*{Question 8}
\textit{Certain stimuli seem to jump off the page at the viewer, causing an involuntary shift of attention that is referred to as:}
\medskip

\begin{tabular}{L{0.3cm} L{14.9cm}}
& A.\ The Stroop effect \\
& \textbf{B.\ Attentional capture \checkmark} \\
& C.\ Attenuation \\
& D.\ The cocktail party effect \\
\end{tabular}

\vspace{0.3cm}
\noindent\textbf{Ans:} B — Attentional capture\\[0.2cm]
\textbf{Why this answer:} \\
\textbf{Attentional capture} refers to the automatic, involuntary redirection of attention toward a salient or conspicuous stimulus in the visual field. Unlike voluntary (endogenous) attention, attentional capture is \textbf{exogenous} — it is driven by the properties of the stimulus itself (e.g., a sudden flash, bright color, or abrupt onset) rather than by the observer's goals or intentions. This concept is central to research on \textbf{visual search} and \textbf{spatial attention}.
\vspace{0.4cm}

\subsubsection*{Question 9}
\textit{Stroop interference lessens when:}
\medskip

\begin{tabular}{L{0.3cm} L{14.9cm}}
& A.\ Participants are better readers \\
& \textbf{B.\ Participants are given more practice at naming colours \checkmark} \\
& C.\ Participants are girls rather than boys \\
& D.\ Participants are encouraged to focus carefully \\
\end{tabular}

\vspace{0.3cm}
\noindent\textbf{Ans:} B — Participants are given more practice at naming colours\\[0.2cm]
\textbf{Why this answer:} \\
The \textbf{Stroop effect} occurs because reading is a highly automatic process that interferes with the controlled process of naming the ink color of a word. \textbf{Practice at colour naming} strengthens the controlled processing pathway, making it more efficient and less susceptible to interference from the automatic word-reading response. Notably, being a better reader would actually \textit{increase} interference (since reading would be even more automatic), while simply being told to focus does not effectively reduce the Stroop effect.
\vspace{0.4cm}

\subsubsection*{Question 10}
\textit{Research on divided attention suggests that:}
\medskip

\begin{tabular}{L{0.3cm} L{14.9cm}}
& A.\ Some people can multitask without any drop in performance \\
& B.\ There are no limits on the number of things that we can successfully do at once \\
& C.\ As individual tasks become more demanding, multitasking becomes more efficient \\
& \textbf{D.\ If you think that you are doing two things simultaneously, you are probably really switching attention back and forth between the two \checkmark} \\
\end{tabular}

\vspace{0.3cm}
\noindent\textbf{Ans:} D — If you think that you are doing two things simultaneously, you are probably really switching attention back and forth between the two\\[0.2cm]
\textbf{Why this answer:} \\
Research on \textbf{divided attention} consistently demonstrates that true simultaneous processing of two demanding tasks is extremely rare. What feels like multitasking is typically \textbf{rapid task-switching} — alternating attention back and forth between tasks. Each switch incurs a \textbf{time cost} and a potential accuracy cost. Performance on both tasks declines as the demands of either individual task increase, reflecting the \textbf{limited capacity} of attentional resources.
\vspace{0.4cm}

\medskip
\subsection*{Assignment Quick Answer Summary}

\begin{center}
\renewcommand{\arraystretch}{1.4}
\begin{tabular}{|L{0.8cm}|L{6.2cm}|L{8.2cm}|}
\hline
\bb{Q\#} & \bb{Topic} & \bb{Correct Answer} \\
\hline
1 & Attenuation Theory & D — Process only as much as necessary \\
2 & Cocktail Party Effect (Shadowing) & B — The listener's name \\
3 & Kahneman's Capacity Model & A — Enduring dispositions \\
4 & Priming & C — Primed \\
5 & Cocktail Party Effect (Real-world) & C — Cocktail party effect \\
6 & Broadbent's Filter Theory & D — Language (semantic, not physical) \\
7 & Automatic vs. Controlled Processing & D — Mental filtering (not a criterion) \\
8 & Attentional Capture & B — Attentional capture \\
9 & Stroop Effect & B — More practice at naming colours \\
10 & Divided Attention & D — Rapid switching, not true simultaneity \\
\hline
\end{tabular}
\end{center}

\medskip
\end{document}
